\section{Magisterial Chronology}\label{app:chronology}

The acts through which the Church has taught and ruled on the angels, in
order of date. ``Locus'' gives the passage on which this study relies;
``Teaching'' states what the act bears on the angels, with the grade the
act supports (see Section~\ref{sec:faith}). Disciplinary acts on the cult
of the angels are included, since they fix the limits of devotion
(Section~\ref{sec:devotion}, Appendix~\ref{app:names}).

\begin{fourstable}{Date}{Act and authority}{Locus}{Teaching (grade)}
325 & Council of Nicaea I (ecumenical), creed & DS 125 & God is maker
``of all things visible and invisible''; the unseen creation belongs to
the one Creator (Defined; read of the angels by the whole tradition)\\
381 & Council of Constantinople I (ecumenical), creed & DS 150 &
\latin{factorem caeli et terrae, visibilium omnium et invisibilium}; the
form professed by the Church (Defined)\\
c.\ 343--381 & Council of Laodicea (particular synod), canon 35 &
can.\ 35 & Forbids Christians to leave the Church and invoke angels in
private assemblies; names the practice ``covert idolatry'' under anathema
(disciplinary; witness to an early abuse)\\
447, 21 July & Leo I, letter \latin{Quam laudabiliter} to Turibius of
Astorga (papal letter) & DS 286 & The devil would be good had he remained
as made; he misused his natural excellence and did not pass into a
contrary substance; evil has no nature (papal teaching; taken up by Braga
and Lateran IV)\\
5th c. & \emph{Statuta Ecclesiae Antiqua} (collection, southern Gaul) &
DS 325 & A bishop-elect is examined on whether ``the devil became evil
not by his condition but by his will'' (disciplinary witness to the
common faith)\\
543 & Edict of Justinian to Menas, published in the Synod of
Constantinople (imperial edict, synodal canons against Origen) &
DS 403--411 & Condemns: souls pre-existing as holy powers (403); the
Word made like each heavenly order (406); ensouled stars as rational
powers (408); Christ crucified for demons in the age to come (409);
temporary punishment and restoration of demons (411) (condemnations;
their contradictories stand as the synod's teaching)\\
561 & Council of Braga I (particular council), anathematisms against the
Priscillianists & DS 455--464 & Condemns: souls or angels from God's
substance (455); souls fallen from heaven into bodies (456); the devil
not first a good angel made by God but from chaos, the principle of evil
(457); the devil maker of creatures, storms, droughts (458); the forming
of bodies and all flesh the work of demons (462--463)\\
745, 25 Oct. & Roman synod under Pope Zachary, against Aldebert
(papal synod) & actio III & Only three angels' names are known from
Scripture --- Michael, Gabriel, Raphael; other names in Aldebert's prayer
are ``rather demons'' invoked; his writings burned (disciplinary)\\
1208, 18 Dec. & Innocent III, \latin{Eius exemplo}, profession for the
Waldensians (papal letter) & DS 797 & ``We believe that the devil became
evil not by his condition but by his will'' (profession of faith
prescribed by the pope)\\
1215, Nov. & Lateran IV (ecumenical), \latin{Firmiter credimus} &
DS 800--801 & God created from nothing, at the beginning of time, both
creatures, the angelic and the earthly; the devil and the other demons
were created good by nature and became evil of themselves; perpetual
punishment with the devil for the reprobate (Defined)\\
1442, 4 Feb. & Council of Florence (ecumenical), \latin{Cantate Domino}
for the Copts & DS 1333, 1336 & All creatures, spiritual and corporeal,
made good by the highest good but mutable because made from nothing; no
nature of evil; two principles condemned (ecumenical teaching confirming
Lateran IV)\\
1546, 17 June & Council of Trent (ecumenical), Session V, original sin &
DS 1511 & Adam's sin brought captivity under the power of the devil,
``who had the empire of death'' (Heb 2:14) (Defined on original sin;
supposes the devil's real power over fallen man)\\
1870, 24 Apr. & Vatican I (ecumenical), \latin{Dei Filius} &
DS 3002, 3025 & Repeats Lateran IV: both creatures made from nothing at
the beginning of time by God's most free counsel; anathema on whoever
denies that all things, spiritual and material, are produced from nothing
as to their whole substance (Defined)\\
1950, 12 Aug. & Pius XII, encyclical \latin{Humani generis} &
n.\ 26; DS 3891 & Lists among reproved opinions the questioning of
whether angels are personal beings (ordinary magisterium; basis for
Church teaching that angels are personal)\\
1968, 30 June & Paul VI, \emph{Credo of the People of God} (solemn papal
profession, expressly not a dogmatic definition) & 3, 8, 29 & Creator
``of things invisible such as the pure spirits which are also called
angels''; the saints associated with the holy angels in Christ's rule
(catechetical restatement of defined doctrine)\\
1972, 15 Nov. & Paul VI, general audience \emph{Liberaci dal male}
(papal catechesis) & DS 800 cited & Denial of the devil's existence goes
outside biblical and ecclesiastical teaching; evil is ``a living being,
spiritual, perverted and perverting''; not every sin is directly
diabolical (catechetical restatement)\\
1975, 26 June & CDF, study \emph{Christian Faith and Demonology},
commended by the Congregation (expert study with the Congregation's
approval) & conclusions & No explicit magisterial affirmation of the
demons' existence, yet each principal point of Lateran's chapter has
equal dogmatic value; the demonic world's existence is a ``dogmatic
datum'' of the ordinary and universal teaching\\
1983, 24 Sept. & CDF, decisions on the Opus Angelorum movement (as
recalled in the decree of 1992) & AAS 76 (1984) 175--176 & Forbade a
cult of the angels using the names drawn from the presumed revelations of
Gabriele Bitterlich (disciplinary; known here only through the 1992
decree)\\
1986, July--Aug. & John Paul II, six general audiences on the angels and
the demons (papal catechesis) & 9, 23, 30 July; 6, 13, 20 Aug. & Restates
the defined doctrine; presents the angels as personal, immortal, pure
spirits; the nine choirs ``without absolute value''; guardian angels;
Satan's fall, his power not infinite, Christ's victory\\
1992, 6 June & CDF, decree on the Opus Angelorum, approved by John Paul
II (congregational decree) & AAS 84 (1992) 805--806 & Theories on the
angels' names and choirs from the movement's presumed revelations are not
taught; consecrations to angels and their ``remote transmission''
prohibited; exorcisms only by the Church's norms (disciplinary)\\
1992; typ.\ 1997 & \emph{Catechism of the Catholic Church} (universal
catechism) & 328--336, 391--395, 414 & The angels' existence ``a truth of
faith''; personal and immortal; created through and for Christ; guardians
of each believer; the fall a free, radical, irrevocable choice; Satan's
power not infinite (Church teaching; what is defined remains defined)\\
2001, 17 Dec. & CDWDS, \emph{Directory on Popular Piety} (congregational
directory) & 213--217 & Devotion to the angels grounded in Scripture and
the feasts of 29 September and 2 October; naming angels beyond Gabriel,
Raphael, and Michael ``should be discouraged'' (disciplinary)\\
2005, 28 June & Benedict XVI, \emph{Compendium} of the Catechism
(catechetical compendium) & 59--61, 74--75 & Angels purely spiritual,
incorporeal, invisible, immortal, personal; their ministry; the fall and
its irrevocability; Satan's power not infinite (Church teaching)\\
\end{fourstable}

Three notes on the table. First, the doctrinal line runs through the
creeds and the anti-dualist acts: Nicaea, Constantinople, Leo I, Braga,
Lateran IV, Florence, Vatican I. What these define is the existence,
creation, goodness, and self-willed fall of the angels. Second, the
disciplinary line runs from Laodicea through the Roman synod of 745 to
the Directory of 2001 and the Opus Angelorum decree: the Church has
repeatedly restricted the cult and naming of angels to what Scripture
gives. Third, the twentieth-century acts --- \latin{Humani generis}, the
\emph{Credo} of Paul VI, the CDF study, the catecheses of John Paul II,
and the Catechism --- neither add definitions nor retract any; they state
which propositions are defined, which are taught, and which remain open.
